\chapter{The learned rollout and the headline statements}

This chapter completes the tier-2 theory with the case that has no rate (T9'),
puts the learned experts into the master error bound, once per coupling tier
(T22), states what the approximation theorem would give if it is assumed (T2),
and ends with the three sentences the proposal can say, each as one theorem
that quotes the others.

\section{T9': non-expansive is enough without a rate}

\paragraph{In plain words.}
A passive coupling gives a sweep that lengthens no distance, but not one that
shortens every distance.  Iterating such a map need not converge: a rotation
never settles.  Averaging each sweep with the current iterate does converge,
to a fixed point of the sweep, whenever the sweep has one.  No rate is claimed.

\begin{definition}[The averaged sweep]\label{def:averaged}
  \lean{Atlas.averaged, Atlas.averaged_fixedPt_iff}
  \leanok
  For $T:E\to E$ and a real $\omega$, $T_\omega(x)=(1-\omega)\,x+\omega\,T(x)$.
  For $\omega\ne0$ it has the same fixed points as $T$.
\end{definition}

\begin{lemma}[T9', the identity]\label{lem:T9p-identity}
  \lean{Atlas.norm_sq_convex_combination}
  \leanok
  In a real inner-product space, for all $a,b$ and all real $\omega$,
  \[
    \norm{(1-\omega)a+\omega b}^2
      =(1-\omega)\norm{a}^2+\omega\norm{b}^2-\omega(1-\omega)\norm{a-b}^2 .
  \]
\end{lemma}

\begin{proof}
  \leanok
  Expand both sides.
\end{proof}

\begin{lemma}[T9', one averaged step]\label{lem:T9p-step}
  \lean{Atlas.averaged_step}
  \leanok
  \uses{def:averaged, lem:T9p-identity}
  Let $\norm{Tx-Ty}\le\norm{x-y}$ for all $x,y$, let $T(p)=p$ and
  $\omega\ge0$.  Then for every $x$,
  \[
    \norm{T_\omega(x)-p}^2\le\norm{x-p}^2-\omega(1-\omega)\norm{x-T(x)}^2 .
  \]
\end{lemma}

\begin{proof}
  \leanok
  \uses{lem:T9p-identity}
  Lemma~\ref{lem:T9p-identity} with $a=x-p$ and $b=T(x)-p$, and
  $\norm{T(x)-p}=\norm{T(x)-T(p)}\le\norm{x-p}$.
\end{proof}

\begin{lemma}[T9', the residuals tend to zero]\label{lem:T9p-residual}
  \lean{Atlas.averaged_residual_tendsto_zero}
  \leanok
  \uses{lem:T9p-step}
  Under the hypotheses of Lemma~\ref{lem:T9p-step} with $0<\omega<1$, the
  iterates $x_{n+1}=T_\omega(x_n)$ satisfy $\norm{x_n-T(x_n)}\to0$.
\end{lemma}

\begin{proof}
  \leanok
  \uses{lem:T9p-step}
  The numbers $\norm{x_n-p}^2$ decrease and are bounded below, so they
  converge, and their successive differences tend to zero; those differences
  bound $\omega(1-\omega)\norm{x_n-T(x_n)}^2$ from above.
\end{proof}

\begin{theorem}[T9', Krasnoselskii--Mann in finite dimensions]\label{thm:T9p}
  \lean{Atlas.krasnoselskii_mann}
  \leanok
  \uses{lem:T9p-step, lem:T9p-residual}
  Let $E$ be a finite-dimensional real inner-product space, $T:E\to E$ with
  $\norm{Tx-Ty}\le\norm{x-y}$ for all $x,y$ and a fixed point, and
  $0<\omega<1$.  For every $x_0$ the iterates $x_{n+1}=T_\omega(x_n)$ converge
  to a fixed point of $T$.
\end{theorem}

\begin{proof}
  \leanok
  \uses{lem:T9p-step, lem:T9p-residual}
  By Lemma~\ref{lem:T9p-step} the iterates stay in a closed ball around a
  fixed point, which is compact, so a subsequence converges to some $q$.  $T$
  is continuous and the residuals tend to zero
  (Lemma~\ref{lem:T9p-residual}), so $T(q)=q$.  By
  Lemma~\ref{lem:T9p-step} with $p=q$, $\norm{x_n-q}$ decreases; it tends to
  zero along the subsequence, hence altogether.
\end{proof}

\section{T22: the master bound with learned experts}

\paragraph{In plain words.}
The master bound (T4) adds up one-step defects.  For a coupled step it names
three: $\tau$, the experts being wrong even with the right interface data;
$\sigma$, the interface problem being posed with the wrong operator; $\gamma$,
the interface solve being stopped early.  At tier 2 the learned experts put
$C_\mu\,\delta/(1-\rho)$ into $\sigma$ and $C_\mu\,\rho^k$ times the initial
interface error into $\gamma$.  At tier 0 there is no iteration, so there is no
$\gamma$, and the learned part of the defect is the Galerkin error that T25
bounds.

\begin{lemma}[T22, tier 2, one step]\label{lem:T22-step}
  \lean{Atlas.tier2_step_defect}
  \leanok
  \uses{thm:T1}
  Let $\Lambda$ be a complete metric space of interface data, and let the step
  $\Phi(\lambda,v)$ satisfy
  $\norm{\Phi(\lambda,v)-\Phi(\lambda',v)}\le C_\mu\,d(\lambda,\lambda')$ at
  the state $v$.  Let $T_c(\lambda^\star)=\lambda^\star$, let $T_l$ be a
  $\rho$-contraction with $0\le\rho<1$, and
  $d(T_l\lambda^\star,T_c\lambda^\star)\le\delta$.  Then $T_l$ has a fixed
  point $\lambda^\dagger$, and for every $\lambda^{(0)}$ and $k$,
  \[
    \norm{\Phi(\lambda^\dagger,v)-\Phi(\lambda^\star,v)}\le\frac{C_\mu\,\delta}{1-\rho},
    \qquad
    \norm{\Phi(T_l^{\,k}\lambda^{(0)},v)-\Phi(\lambda^\dagger,v)}
      \le C_\mu\,\rho^k\,d(\lambda^{(0)},\lambda^\dagger).
  \]
\end{lemma}

\begin{proof}
  \leanok
  \uses{thm:T1}
  Theorem~\ref{thm:T1} gives $\lambda^\dagger$, its distance to
  $\lambda^\star$ and the rate; then the Lipschitz bound of the step.
\end{proof}

\begin{theorem}[T22, tier 2]\label{thm:T22-tier2}
  \lean{Atlas.learned_master_bound_tier2}
  \leanok
  \uses{lem:T22-step, cor:T4-three}
  Let the step use the datum returned by $k$ learned sweeps,
  $u^{n+1}=\Phi\bigl(T_l(u^n)^k\lambda^{(0)}(u^n),\,u^n\bigr)$, and magnify
  the difference between the computed and the true trajectory by at most
  $L\ge0$.  Along the true trajectory $u^{n,\star}$ let the hypotheses of
  Lemma~\ref{lem:T22-step} hold with constants $C_\mu,\rho,\delta$, and let
  the starting guess be within $D$ of the learned fixed point.  Then
  \[
    \norm{u^N-u^{N,\star}}\le L^N\norm{u^0-u^{0,\star}}
      +\sum_{n=1}^{N}L^{N-n}\Bigl(\norm{\tau^n}
        +\frac{C_\mu\,\delta}{1-\rho}+C_\mu\,\rho^k\,D\Bigr),
  \]
  with $\tau^n=\Phi(\lambda^{n-1,\star},u^{n-1,\star})-u^{n,\star}$.
\end{theorem}

\begin{proof}
  \leanok
  \uses{lem:T22-step, cor:T4-three}
  T4 with the three terms named, and Lemma~\ref{lem:T22-step} at every step.
\end{proof}

\begin{theorem}[T22, tier 0]\label{thm:T22-tier0}
  \lean{Atlas.learned_master_bound_tier0, Atlas.learned_master_bound_tier0_superelement}
  \leanok
  \uses{lem:T4-accumulation, lem:energy-triangle, thm:T25-bound}
  In the energy size $\norm{\cdot}_m$ of a quadratic energy: let
  $u^{n,\mathrm{ex}}$ be the undivided discrete step from the true state
  $u^{n,\star}$, and $\tilde u^n$ the tier-0 step from it.  If
  $\norm{u^{n+1}-\tilde u^n}_m\le L\norm{u^n-u^{n,\star}}_m$,
  $\norm{u^{n,\mathrm{ex}}-u^{n+1,\star}}_m\le\tau^n$ and
  $\norm{u^{n,\mathrm{ex}}-\tilde u^n}_m\le g^n$, then
  \[
    \norm{u^N-u^{N,\star}}_m\le L^N\norm{u^0-u^{0,\star}}_m
      +\sum_{n=0}^{N-1}L^{N-1-n}\bigl(\tau^n+g^n\bigr).
  \]
  There is no incomplete-solve term.  For the tier-0 solve of
  Definition~\ref{def:superelement}, under assumptions 1 and 2, $g^n$ may be
  taken as the bound of Theorem~\ref{thm:T25-bound}: the error of any reference
  field plus the energy of the network's field errors.
\end{theorem}

\begin{proof}
  \leanok
  \uses{lem:T4-accumulation, lem:energy-triangle, thm:T25-bound}
  $u^{n+1}-u^{n+1,\star}=(u^{n+1}-\tilde u^n)+(\tilde u^n-u^{n,\mathrm{ex}})
  +(u^{n,\mathrm{ex}}-u^{n+1,\star})$, the triangle inequality of the energy
  size, and the discrete Gronwall inequality.
\end{proof}

\section{T2: existence, as a conditional theorem}

\paragraph{In plain words.}
Can a learned coupling do as well as the classical one?  If a network can
approximate the classical interface map and its derivative to within $\eta$ on
a ball around the classical answer, then its iteration contracts at the
classical rate plus $\eta$ and its answer is within $\eta/(1-\rho-\eta)$ of
the classical one.  The ``if'' is the universal approximation theorem in
$C^1$, which is cited and not proved here.  Whether a trained network meets
it is measured, one network at a time.

\begin{theorem}[T2, conditional]\label{thm:T2}
  \lean{Atlas.existence_conditional, Atlas.existence_conditional_c1,
    Atlas.existence_of_learned_coupling}
  \leanok
  \uses{thm:T1}
  Let $E$ be a real Banach space, $B$ the closed ball of radius $r\ge0$
  around $a_c$, $T_c(a_c)=a_c$ and
  $\norm{T_cx-T_cy}\le\rho\norm{x-y}$ on $B$.  Let $T_l$ satisfy
  $\norm{T_la_c-T_ca_c}\le\eta$ and
  $\norm{(T_l-T_c)x-(T_l-T_c)y}\le\eta\norm{x-y}$ on $B$ (which holds if
  $T_l-T_c$ has a derivative of norm at most $\eta$ on $B$).  If
  $\rho+\eta<1$ and $\eta\le(1-\rho-\eta)\,r$, then $T_l$ has exactly one
  fixed point $a_l$ in $B$; from every $x_0\in B$ the iterates stay in $B$ and
  $\norm{T_l^{\,k}x_0-a_l}\le(\rho+\eta)^k\norm{x_0-a_l}$; and
  \[
    \norm{a_l-a_c}\le\frac{\eta}{1-\rho-\eta}.
  \]
  Consequently, if a class of maps contains one such $T_l$, it contains a map
  with these three properties.
\end{theorem}

\begin{proof}
  \leanok
  \uses{thm:T1}
  $T_l=T_c+(T_l-T_c)$ is $(\rho+\eta)$-Lipschitz on $B$, and
  $\norm{T_lx-a_c}\le(\rho+\eta)r+\eta\le r$, so $T_l$ maps $B$ into itself.
  $B$ is complete; Theorem~\ref{thm:T1} on $B$.  The derivative form is the
  mean value inequality on the convex set $B$.
\end{proof}

\section{The headline statements}

\begin{theorem}[Headline, tier 0]\label{thm:headline-tier0}
  \lean{Atlas.tier0_headline}
  \leanok
  \uses{thm:T25-optimal, thm:T25-solvable}
  \emph{Whatever a learned superelement returns, the coupled system is
  solvable, and its answer is the best its modes can represent in the energy
  norm.}  Under the three assumptions of Definition~\ref{def:assumptions},
  with independent port patterns and every free coordinate read by some piece:
  for any fields the network returned, $\tilde{\mathsf S}c=\tilde\chi$ has
  exactly one solution $c$, and
  $\norm{u-w(c)}_M\le\norm{u-w(c')}_M$ for every $c'$.
\end{theorem}

\begin{proof}
  \leanok
  \uses{thm:T25-optimal, thm:T25-solvable}
  Theorems~\ref{thm:T25-solvable} and~\ref{thm:T25-optimal}.
\end{proof}

\begin{theorem}[Headline, tier 2]\label{thm:headline-tier2}
  \lean{Atlas.tier2_headline}
  \leanok
  \uses{thm:T9a, thm:T22-tier2}
  \emph{If every learned expert's scattering map is certified to contract by
  $\rho<1$ and differs from the classical one by at most $\delta$, the coupled
  iteration converges geometrically from any start to a point within
  $\delta/(1-\rho)$ of the classical answer; over $N$ steps the error obeys
  the master bound with that term added.}  Here the interface sweep of a step
  is the wave sweep of Definition~\ref{def:wave-sweep}, and
  $\delta=\bigl(\sum_i\delta_i^2\bigr)^{1/2}$ combines the experts' defects at
  the classical answer.
\end{theorem}

\begin{proof}
  \leanok
  \uses{thm:T9a, thm:T22-tier2}
  Theorem~\ref{thm:T9a} at every step, and Theorem~\ref{thm:T22-tier2} with
  the wave sweeps as the interface maps.
\end{proof}

\begin{theorem}[Headline, defect correction]\label{thm:headline-dc}
  \lean{Atlas.defect_correction_headline}
  \leanok
  \uses{thm:T5}
  \emph{Inside defect correction, a learned expert can change how many
  classical steps the answer takes.  It cannot change the answer.}  Run defect
  correction with two cheap maps $\Psi_1,\Psi_2$.  If both runs converge, to
  $a_1$ and $a_2$, and the maps are continuous there, then
  $\Phi(a_1)=a_1$ and $\Phi(a_2)=a_2$; if $\Phi$ has at most one fixed point,
  $a_1=a_2$.
\end{theorem}

\begin{proof}
  \leanok
  \uses{thm:T5}
  Theorem 1 of T5, twice.
\end{proof}
